\Cref{tab:1-contributions}列出三项关联项目贡献及继承基础。
核心是具体表示--执行--预算组合的适配与评价。
JHQ、Cartesian量化、IVF、bitonic／radix、merge与CUDA Graph仍属于前人工作。

\par\medskip\noindent\begin{minipage}{\linewidth}
\centering\captionsetup{hypcap=false}
\captionof{table}[技术贡献与证据边界]{技术贡献与证据边界。}
\label{tab:1-contributions}
\footnotesize\setlength{\tabcolsep}{4pt}
\begin{tabularx}{\linewidth}{@{}l*{3}{>{\raggedright\arraybackslash}X}@{}}
\toprule
 & 继承基础 & 项目贡献 & 证据 \\
\midrule
C1 & JHQ旋转、码本与层次 & GPU训练／编码、有界batch构建及列表有序packed表示
& 17.78M向量搜索、常驻20.44\,GiB、重复train+add 10.54--10.65\,s、采样构建峰值22.58\,GiB \\
C2 & Cartesian可加性与既有GPU原语 & 融合拆表评分、完整键分组保留／界、lane精排与调度
& 两组LUT在22/24单元领先；batch 8调度$2.727\times$；明确正确性范围 \\
C3 & 候选扩张和输出一致性思路 & 工作负载预算规则、有条件选择分析、留出对照与成本／风险评估
& FULL的12/12单元相对损失达标；小样本揭示风险；四个单元直接测量校准复用 \\
\bottomrule
\end{tabularx}
\end{minipage}\par\medskip

E1大batch的0.97附近插值比较中，JHQ在三个数据集
领先，arXiv与RaBitQ接近平手，其他recall区间有对手领先。
FULL最大实测留出损失相对最佳被测臂为0.0006，$S=128$一个单元有40/64次越阈。
这些是经验质量结果；10个FULL单元一般嵌套仍有条件，但十二个FULL有限网格选择均与枚举一致（\cref{sec:7-budget-trace}）。
OpenAI E1使用TF32 GEMM；独立FP32敏感性仍保留较小领先（\cref{sec:7-precision}）。

本COMP8800项目由Rui Luo在Dr Mengxuan Zhang指导下呈现。
ADC合著稿列Rui Luo与Mengxuan Zhang，复用JHQ-GPU实现及测量产物\cite{adc_artifact}。
论文扩展正确性契约、新版实验分析与失败方案讨论，历史稿数字保留各自版本。
技术贡献标签界定项目范围；来源说明区分继承与共同工作，并披露AI协助。
论文作者对提交的主张与归属标注负责。
